On applique la formule~: $W_{AB}(\vec{P})=mg(z_{A}-z_{B})$.
\begin{enumerate}
    \item Entre le début et la fin du couloir, $z_{A}=z_{B}$ puisque le couloir
          est horizontal~; donc~:
          \[
              W_1(\vec{P})=0
          \]
    
          Il revient au même de dire que le poids $\vec{P}$ est continuellement
          perpendiculaire au déplacement de son point d'application.

    \item Entre le rez-de-chaussée et le troisième étage, le point d'application
          du poids s'élève de $h=3 H=3 \times 2,80=\qty{8.40}{\m}$. Donc~:
          \begin{itemize}
              \item le travail du poids est négatif
              \item $W_2(\vec{P})=-mgh$.
          \end{itemize}
    
          Application numérique~:
          \[
              W_2(\vec{P})=-8,0 \times 10 \times 8,40 \approx \qty{-6.7e2}{\J}.
          \]
\end{enumerate}

